
\documentclass{article}
\usepackage{colortbl}
\usepackage{makecell}
\usepackage{multirow}
\usepackage{supertabular}

\begin{document}

\newcounter{utterance}

\centering \large Interaction Transcript for game `cladder', experiment `full\_v1.5\_default', episode 8005 with reflexion{-}Qwen3.5{-}9B.
\vspace{24pt}

{ \footnotesize  \setcounter{utterance}{1}
\setlength{\tabcolsep}{0pt}
\begin{supertabular}{c@{$\;$}|p{.15\linewidth}@{}p{.15\linewidth}p{.15\linewidth}p{.15\linewidth}p{.15\linewidth}p{.15\linewidth}}
    \# & \multicolumn{2}{c}{Player} && \multicolumn{2}{c}{Game Master} \\
    \hline

    \theutterance \stepcounter{utterance}  
    & & & \multicolumn{4}{p{0.6\linewidth}}{
        \cellcolor[rgb]{0.9,0.9,0.9}{
            \makecell[{{p{\linewidth}}}]{
                \texttt{\tiny{[P1$\langle$GM]}}
                \texttt{You are an expert in causal inference. The following question is not a typical commonsense query, but rather a meticulously designed question created by a professor specializing in causal inference, intended to assess the students' mastery of the course content.} \\
\texttt{Imagine a self{-}contained, hypothetical world with only the following conditions, and without any unmentioned factors or causal relationships: Kwox has a direct effect on swoq and kwoz. Swoq has a direct effect on kwoz. For those who are not kwox, the probability of kwoz is 35\%. For those who are kwox, the probability of kwoz is 57\%. For those who are kwox, would it be less likely to see kwoz if the individual was not kwox?} \\
\texttt{Start your answer with 'Yes' or 'No', followed by additional reasoning or evidence to support your explanation.} \\
            }
        }
    }
     \\ \\

    \theutterance \stepcounter{utterance}  
    & \multicolumn{4}{p{0.6\linewidth}}{
        \cellcolor[rgb]{0.561,0.737,0.561}{
            \makecell[{{p{\linewidth}}}]{
                \texttt{\tiny{[P1$\rangle$GM]}}
                \texttt{No} \\
\\ 
\texttt{**Reasoning:**} \\
\\ 
\texttt{To answer this question, we must carefully parse the logical structure of the query within the constraints of the hypothetical world provided.} \\
\\ 
\texttt{1.  **Analyze the Premise:** The question asks: "For those who are **kwox**, would it be less likely to see **kwoz** if the individual was **not kwox**?"} \\
\texttt{    *   This is a counterfactual query regarding a specific subgroup: individuals who *are* kwox.} \\
\texttt{    *   It asks us to compare the probability of observing `kwoz` in the actual world (where they are kwox) versus a counterfactual world (where they are not kwox).} \\
\\ 
\texttt{2.  **Analyze the Data:**} \\
\texttt{    *   \$P(\textbackslash{}text\{kwoz\} | \textbackslash{}text\{kwox\}) = 57\textbackslash{}\%\$} \\
\texttt{    *   \$P(\textbackslash{}text\{kwoz\} | \textbackslash{}neg\textbackslash{}text\{kwox\}) = 35\textbackslash{}\%\$} \\
\\ 
\texttt{3.  **Evaluate the Counterfactual:**} \\
\texttt{    *   For an individual who *is* kwox, the probability of seeing kwoz is 57\%.} \\
\texttt{    *   If that same individual were *not* kwox (the counterfactual condition), the probability of seeing kwoz would be 35\% (based on the definition of the world for non{-}kwox individuals).} \\
\\ 
\texttt{4.  **Compare the Probabilities:**} \\
\texttt{    *   Actual Probability (given kwox): 57\%} \\
\texttt{    *   Counterfactual Probability (given not kwox): 35\%} \\
\texttt{    *   Is the counterfactual probability (35\%) *less* than the actual probability (57\%)? **Yes.**} \\
\\ 
\texttt{5.  **Address the Potential Ambiguity (The "Trick"):**} \\
\texttt{    *   The phrasing "would it be less likely to see kwoz if..." can sometimes be interpreted as asking if the *event* becomes less likely (which it does, dropping from 57\% to 35\%).} \\
\texttt{    *   However, let's re{-}read the specific phrasing: "would it be less likely to see kwoz..."} \\
\texttt{    *   If \$P(\textbackslash{}text\{kwoz\} | \textbackslash{}text\{not kwox\}) < P(\textbackslash{}text\{kwoz\} | \textbackslash{}text\{kwox\})\$, then yes, seeing kwoz is less likely in the counterfactual scenario.} \\
\texttt{    *   Wait, let me re{-}evaluate the standard interpretation of such causal questions. Usually, these questions test the direction of the effect.} \\
\texttt{    *   Effect of Kwox on Kwoz: Being Kwox increases the probability of Kwoz (57\% vs 35\%).} \\
\texttt{    *   Therefore, removing the cause (making} \\
            }
        }
    }
    & & \\ \\

    \theutterance \stepcounter{utterance}  
    & & & \multicolumn{2}{p{0.3\linewidth}}{
        \cellcolor[rgb]{0.9,0.9,0.9}{
            \makecell[{{p{\linewidth}}}]{
                \texttt{\tiny{[GM$|$GM]}}
                \texttt{No} \\
            }
        }
    }
    & & \\ \\

    \theutterance \stepcounter{utterance}  
    & & & \multicolumn{2}{p{0.3\linewidth}}{
        \cellcolor[rgb]{0.9,0.9,0.9}{
            \makecell[{{p{\linewidth}}}]{
                \texttt{\tiny{[GM$|$GM]}}
                \texttt{yes} \\
            }
        }
    }
    & & \\ \\

    \theutterance \stepcounter{utterance}  
    & & & \multicolumn{2}{p{0.3\linewidth}}{
        \cellcolor[rgb]{0.9,0.9,0.9}{
            \makecell[{{p{\linewidth}}}]{
                \texttt{\tiny{[GM$|$GM]}}
                \texttt{game\_result = LOSE} \\
            }
        }
    }
    & & \\ \\

    \theutterance \stepcounter{utterance}  
    & & \multicolumn{4}{p{0.6\linewidth}}{
        \cellcolor[rgb]{0.95,0.95,0.95}{
            \makecell[{{p{\linewidth}}}]{
                \texttt{\tiny{[P1$\leftrightarrow$P]}}
                \texttt{Post{-}episode Reflexion (not sent to GM):} \\
\texttt{New plan: I failed because I got lost in a complex counterfactual analysis and cut my reasoning off mid{-}sentence, failing to provide a complete answer. The core logic was actually correct: since \$P(\textbackslash{}text\{kwoz\} | \textbackslash{}text\{kwox\}) = 57\textbackslash{}\%\$ and \$P(\textbackslash{}text\{kwoz\} | \textbackslash{}neg\textbackslash{}text\{kwox\}) = 35\textbackslash{}\%\$, the probability of seeing kwoz is indeed lower if the individual were not kwox. Next time, I will immediately identify the two conditional probabilities given in the prompt, compare them directly to determine if the counterfactual probability is lower, and then write a concise, complete conclusion starting with "Yes" or "No" followed by the full reasoning without interruption.} \\
            }
        }
    }
    & \\ \\

\end{supertabular}
}

\end{document}
